% Generated draft: requires mathematical and source review.
\documentclass[11pt]{article}
\usepackage[margin=25mm]{geometry}
\usepackage{fontspec}
\setmainfont{Latin Modern Roman}
\usepackage{amsmath,amssymb,mathtools,unicode-math}
\setmathfont{Latin Modern Math}
\usepackage{enumitem,needspace,xurl,hyperref}
\setlength{\parindent}{0pt}
\setlength{\parskip}{5pt}
\setlength{\emergencystretch}{3em}
\newcommand{\N}{\mathbb{N}}
\newcommand{\Z}{\mathbb{Z}}
\newcommand{\Q}{\mathbb{Q}}
\newcommand{\R}{\mathbb{R}}
\newcommand{\C}{\mathbb{C}}
\newcommand{\F}{\mathbb{F}}
\newcommand{\A}{\mathbb{A}}
\newcommand{\PP}{\mathbb{P}}
\newcommand{\E}{\mathbb{E}}
\newcommand{\eps}{\varepsilon}
\newcommand{\dd}{\,\mathrm{d}}
\begin{document}
\section*{UnsolvedMath ID: 30005501}
\textbf{Source problem code:} catalogue entry
\subsection{Definitions and mathematical statement}
The supplied question concerns blocks of finite group algebras, but does not specify the coefficient field. The following formulation uses the conventional algebraically closed modular setting; this field hypothesis is an explicit interpretive assumption requiring source confirmation.

Let \(p\) be a prime, let \(k\) be an algebraically closed field of characteristic \(p\), and let \(G\) be a finite group. A block algebra of \(kG\) is \(B=kGe\), where \(e\) is a primitive central idempotent of \(kG\). For a \(p\)-subgroup \(P\leq G\), write
\[
\operatorname{Br}_P:(kG)^P\longrightarrow kC_G(P),\qquad
\operatorname{Br}_P\!\left(\sum_{g\in G}a_g g\right)=\sum_{g\in C_G(P)}a_g g,
\]
where \((kG)^P\) denotes the subalgebra fixed by conjugation by \(P\), and \(C_G(P)\) is its centralizer in \(G\). A defect group of \(B\) is a \(p\)-subgroup \(D\leq G\) maximal under inclusion subject to \(\operatorname{Br}_D(e)\neq0\). Nontrivial defect means \(D\neq\{1\}\).

Define the first Hochschild cohomology over \(k\) by
\[
HH^1_k(B)=\operatorname{Der}_k(B)/\operatorname{Inn}(B),
\]
where \(\operatorname{Der}_k(B)\) consists of the \(k\)-linear maps \(\delta:B\to B\) satisfying \(\delta(xy)=\delta(x)y+x\delta(y)\), and \(\operatorname{Inn}(B)\) consists of maps \(x\mapsto ax-xa\) for \(a\in B\).

\textbf{Question, under the coefficient-field convention above.} For every prime \(p\), every algebraically closed field \(k\) of characteristic \(p\), every finite group \(G\), and every block algebra \(B\) of \(kG\) with nontrivial defect group, must
\[
HH^1_k(B)\neq0
\]
hold?

The original extract refers to ``[9, Question 7.4]'', whose bibliographic identity is not supplied. It also appends a corrupted formula introduced by \(p\geq3\), containing an apparent summation from \(j=0\) to \(w-1\), \(\dim_k HH^1(B)\), and a term rendered as \(\texttt{dimFp (Z(Bpj ))}\). Neither the notation nor the hypotheses of that formula can be recovered reliably from the supplied record. It is therefore retained as an unresolved source qualification, not asserted as an identity. The supplied clean question places no restriction \(p\geq3\) on the nonvanishing question itself.
\subsection{Short English statement}
Does every block of a finite group algebra with nontrivial defect admit a \(k\)-linear derivation that is not inner? The displayed formulation assumes an algebraically closed field of positive characteristic, a convention absent from the supplied statement and requiring confirmation. The record labels the question open; this is supplied curation rather than an independently verified status.
\subsection{Sources}
\begin{itemize}
\item Local UnsolvedMath registry: supplied complete record with permanent ID \(30005501\), title ``First Hochschild Cohomology of Nontrivial-Defect Blocks'', and exact source problem code \texttt{catalogue entry}. The record attributes the problem to the Oberwolfach Reports collection and gives proposed year \(2023\); no individual proposer is named.
\item \emph{Representations of Finite Groups} (2023), Oberwolfach Reports, pp.~1031--1084, DOI \texttt{10.4171/owr/2023/19}: \url{https://doi.org/10.4171/owr/2023/19}. This is the supplied authoritative report citation. The original extract cites ``[9, Question 7.4]'', but the corresponding bibliography entry and precise problem location are not supplied.
\item Supplied literature reference: Benson--Kessar--Linckelmann, \emph{Hochschild cohomology II}: \url{https://arxiv.org/abs/2301.03185}.
\item Supplied literature reference: \emph{Nontriviality of first Hochschild cohomology}: \url{https://doi.org/10.1016/j.jpaa.2022.107192}.
\item Supplied literature reference: \emph{Oberwolfach Report 19/2023}: \url{https://doi.org/10.4171/owr/2023/19}.
\item Literature assessment dated exactly \texttt{2026-08-22}: the supplied curation labels the general assertion open and reports proofs for symmetric-group blocks and several additional families, without specifying those additional families. This is supplied dated curation derived from the OWR collection, not an independently verified literature assessment; its status and cited claims require independent checking.
\item The registry assigns statement status \texttt{corrected\_verified} and states that its reconstruction was checked against the official Oberwolfach report. This is a supplied verification claim, not verification performed for this draft.
\end{itemize}
\end{document}
